\begin{proposition}[Proved: nil-invariance in uniquely divisible directions]
\label{C:divisible_nil_invariance}
Let $k$ be any commutative ring, $G$ a finite group of order $N$,
and $J\subset k$ a nil ideal such that multiplication by $N$ is a
bijection on the additive group of $J$. Then reduction induces
\[
 R_k(G)\xrightarrow{\ \sim\ }R_{k/J}(G).                \tag{C.1.17}
\]
It respects internal tensor products and exterior operations.
In particular, if the additive group of $J$ is a $\mathbb Q$-vector
space, reduction gives an isomorphism of the entire induction rings
$R_k(\mathfrak S_*)\simeq R_{k/J}(\mathfrak S_*)$, and
$\vartheta_k$ is an isomorphism if and only if $\vartheta_{k/J}$ is.
There is no noetherian or finite-generation hypothesis on $J$.
\end{proposition}
\begin{proof}
First suppose $J^2=0$. Finite projectives lift across $J$: represent
one by a finite idempotent matrix, lift the matrix, and replace an
approximate idempotent $e$ by $3e^2-2e^3$; its error is a multiple
of $(e^2-e)^2$ and hence zero. This is the finite case of
\cite[Tags 07LV and 0AMF, Lemma 10.32.7]{C-StacksNil}.
Maps between such projectives lift by projectivity, and any lift
of an isomorphism is invertible, by a finite geometric series in
the square-zero error.

Fix a projective lift $P$ of a representation $\overline P$ over
$k/J$. Lift each of the finitely many action matrices to an
automorphism $r_g$ of $P$, taking $r_1=1$. Let
$D=\operatorname{Hom}_k(P,JP)$; it is square-zero as an ideal in
$\operatorname{End}_k(P)$, and $N$ acts invertibly on $D$ because
$P$ is a direct summand of a finite free module. Conjugation by
the $r_g$ defines an actual $G$-action on $D$: all group-law
errors lie in $D$, whose conjugation action on itself is trivial.
Write
\[
 r_gr_hr_{gh}^{-1}=1+c(g,h),\qquad c(g,h)\in D.
\]
Associativity gives
\[
 c(g,h)+c(gh,l)=g\,c(h,l)+c(g,hl).
\]
Define $a(g)=N^{-1}\sum_{l\in G}c(g,l)$, where division takes
place in $D$; it does not require $N$ to be a unit in $k$.
Summing the preceding identity over $l$ gives
\[
 c(g,h)=a(g)+g\,a(h)-a(gh).
\]
Consequently $r'_g=(1-a(g))r_g$ obeys the group law. Every
representation over $k/J$ has therefore lifted to one over $k$.

Next let $P,Q$ already carry lifted actions, and let
$\bar f:\overline P\to\overline Q$ be equivariant. Choose an
underlying $k$-linear lift $f$. The discrepancies
\[
 d(g)=g\,f-f\in\operatorname{Hom}_k(P,JQ)
\]
obey $d(gh)=d(g)+g\,d(h)$. For
$b=N^{-1}\sum_h d(h)$, summing this identity gives
$d(g)=b-g\,b$. Hence $f+b$ is an equivariant lift. In particular
an equivariant isomorphism lifts to one. This proves uniqueness
of lifts up to isomorphism as well as surjectivity on all Hom sets;
it does not assert that reduction is faithful on Hom sets.

Admissible relations lift too. Given
$0\to\overline A\to\overline E\to\overline B\to0$, lift its
three objects and lift the equivariant surjection to $p:E\to B$.
An underlying $k$-linear lift $s:B\to E$ of its chosen splitting
has $ps$ invertible, so replacing $s$ by $s(ps)^{-1}$ makes
$p$ a split surjection. Its kernel is projective over $k$ and
reduces to $\overline A$; the uniqueness just proved identifies
it equivariantly with the chosen lift $A$. Thus both the objects
and all defining relations of the two Grothendieck groups
correspond, proving (C.1.17) when $J^2=0$.

If $J$ is nilpotent, apply this argument successively to
$J^r/J^{r+1}$. Multiplication by $N$ is bijective on every power
$J^r$: surjectivity follows by dividing one factor of each product,
and injectivity follows from $J^r\subset J$. It is consequently
bijective on each quotient layer. The square-zero argument applies
at every stage.

Finally let $J$ merely be nil. The inverse of multiplication by
$N$ is $k$-linear on $J$. For a finite set $j_1,\ldots,j_t\in J$,
let $L$ be the ideal it generates, and put
\[
 L^{[N]}=\sum_{a\geq0}N^{-a}L\ \subset J.
\]
This is an ideal on which $N$ acts bijectively. It is nilpotent:
the finitely generated nil ideal $L$ has some $L^r=0$, and a
product of $r$ elements of $L^{[N]}$, after multiplication by a
suitable power of $N$, is a sum of products from $L^r$.
Injectivity of that power on $J$ forces the product to be zero.
The ideals $L^{[N]}$ form a filtered family with union $J$.
Thus $k/J=\operatorname*{colim}_L k/L^{[N]}$.
Proposition~\ref{C:filtered_colimits}, whose finite-matrix descent
also includes equivariant maps and admissible sequences, now
extends (C.1.17) from every $L^{[N]}$ to $J$.

Reduction commutes with tensor products, exterior powers and
induction. If $J$ is a $\mathbb Q$-vector space the hypothesis
holds for every $N$, so these degreewise isomorphisms give the
last assertion and carry the polynomial generators to themselves.
\end{proof}

\begin{corollary}[Proved: nonreduced mixed-characteristic positive examples]
\label{C:divisible_nil_positive}
Let $A$ be any ring satisfying C.1 and let $J$ be a nil ideal in
a commutative extension $k$ with $k/J\simeq A$. If $J$ is uniquely
divisible as an additive group, then $k$ satisfies C.1.
For example
\[
 k=\mathbb Z_{(2)}\oplus\mathbb Q\epsilon,
 \qquad\epsilon^2=0,\qquad
 (a,x)(b,y)=(ab,ay+bx)                                \tag{C.1.18}
\]
is local, nonreduced and nonnoetherian, and
$\vartheta_k$ is an isomorphism in every degree. Its closed
fibre is $\mathbb F_2$, its generic fibre is
$\mathbb Q[\epsilon]/\epsilon^2$, and $K_0(k)=\mathbb Z$.
\end{corollary}
\begin{proof}
Apply Proposition~\ref{C:divisible_nil_invariance}.
In (C.1.18) the maximal ideal is
$2\mathbb Z_{(2)}\oplus\mathbb Q\epsilon$; reduction by the
square-zero rational ideal is $\mathbb Z_{(2)}$, covered by
Proposition~\ref{C:DVR}. The ideal $\mathbb Q\epsilon$ is not
finitely generated over $k$: a finite collection has bounded
$2$-denominators, and multiplying by $k$ cannot enlarge them.
This proves nonnoetherianity. The fibre statements follow by
tensoring with $\mathbb F_2$ and $\mathbb Q$, respectively, and
finite projectives over the local ring are free.
\end{proof}

\begin{corollary}[Proved: ordinary fibre and projective-class data are insufficient]
\label{C:fibre_invisible_kernel}
For $A=\mathbb Z_{(2)}$ the maps
$A\to A[\epsilon]/\epsilon^2\to A$ induce the same underlying
spectrum and residue fields. On every integer localization, and
on every base change to a field, their ordinary $K_0$ groups are $\mathbb Z$
and these maps induce the identity. Nevertheless the first ring
satisfies C.1 and the second has the infinite kernel (C.1.12).
Every element of that kernel is zero on restriction to each
residue-field point of $\operatorname{Spec}k$. Thus ordinary
$K_0$ data and these pointwise representation classes alone cannot
give a necessary and sufficient test for C.1. This does not rule
out a criterion that also uses further structure of the ring $k$.
\end{corollary}
\begin{proof}
The underlying spectrum is unchanged by a nilpotent ideal.
The local rings in question, their integer localizations, and
their base changes to any residue field are local; hence their
$K_0$ groups are $\mathbb Z$, with the unit classes matched.
The positive assertion for $A$ is Proposition~\ref{C:DVR}, while
Theorem~\ref{C:dualnumber_C2_ring} gives the full kernel for the
dual-number ring. At its characteristic-zero point each $\eta_q$
vanishes generically. At its characteristic-two point
$\epsilon=0$ and both signs become trivial, so the two terms
of (C.1.11) are the same $2\deg(q)$ trivial representation.
This proves vanishing at every field point without discarding
the nonzero integral classes.
\end{proof}

The finite certificate \texttt{code/C\_finish\_nilfibres.py}
checks the two averaging homotopy identities before division by
the group order on every basis cochain for $C_2,C_3,\mathfrak S_3$
acting on the endomorphisms of their indicated permutation modules:
$89$ basis vectors in degree one and $421$ in degree two. These
integer identities verify the signs used to lift actions and maps;
the arbitrary-ring extension is the proof above.
